\documentclass[12pt,a4paper]{ctexart}
% 公共排版文件（使用 XeLaTeX 编译）
\usepackage[a4paper,margin=2.35cm]{geometry}
\usepackage{booktabs,longtable,tabularx,array,multirow,enumitem}
\usepackage{graphicx,xcolor,amsmath,hyperref,lastpage,fancyhdr}
\usepackage{tikz}
\usetikzlibrary{arrows.meta,positioning,shapes.geometric,fit,calc}
\hypersetup{colorlinks=true,linkcolor=blue!60!black,urlcolor=blue!60!black,citecolor=blue!60!black,pdftitle={逸仙活动云课程报告}}
\setlist{itemsep=2pt,topsep=3pt,leftmargin=2em}
\setlength{\parindent}{2em}
\setlength{\headheight}{15pt}
\renewcommand{\arraystretch}{1.35}
\newcolumntype{Y}{>{\raggedright\arraybackslash}X}
\newcommand{\reportcover}[1]{%
  \begin{titlepage}
  \begin{center}
  \vspace*{3cm}{\zihao{1}\bfseries 中山大学\\[0.6cm]《软件工程》期末大作业\par}
  \vspace{2.5cm}{\zihao{2}\bfseries #1\par}
  \end{center}
  \vspace{2cm}
  \begin{center}
  \begin{tabular}{l}
  {\Large 项目名称：逸仙活动云——中山大学校园活动云平台}\\
  {\Large 小组成员：钟梓俊\quad 邱剑波\quad 张宸睿}\\
  {\Large 指导老师：王可泽}
  \end{tabular}
  \end{center}
  \vfill
  \begin{center}
  {\large 2026 年 7 月 16 日\par}
  \end{center}
  \end{titlepage}}

\pagestyle{fancy}\fancyhf{}\lhead{逸仙活动云}\rhead{软件工程期末大作业}\cfoot{第 \thepage\ 页，共 \pageref{LastPage} 页}
\setcounter{tocdepth}{2}

\begin{document}
\reportcover{软件工程化说明文档}
\pagenumbering{Roman}\tableofcontents\clearpage\pagenumbering{arabic}
\docmeta

\section{软件工程化目标}
本项目以“每次变更可追踪、可构建、可验证、可交接”为工程化目标，围绕软件过程、协作、自动化和知识沉淀建立轻量但可执行的机制。其核心不是堆砌工具，而是让需求—设计—代码—测试—发布形成闭环，并使前后端能在接口演化中保持一致。具体而言，工程化实践旨在达成以下目标：
\begin{enumerate}
\item 让每个功能从需求提出到上线发布的过程可追溯，变更历史可审计。
\item 让构建和测试过程自动化执行，减少人工回归成本，及早发现集成问题。
\item 让环境配置和部署流程可复现，消除“在我机器上能跑”的侥幸心理。
\item 让知识（设计决策、接口契约、部署步骤）以文档形式沉淀，不依赖口头传递。
\end{enumerate}

\section{项目协作方式}
\subsection{团队分工}
项目采用前后端分离的分工模式，团队成员按技术栈分为两个小组：
\begin{enumerate}
\item \textbf{前端组：} 负责 Vue 3 单页应用的开发，包括页面组件、路由配置、状态管理、API 封装和接口适配。工作产出包括前端源代码、页面设计文档和接口联调验证。
\item \textbf{后端组：} 负责 Flask REST API 的开发，包括业务逻辑、数据模型、认证授权、外部服务集成和部署配置。工作产出包括后端源代码、数据库模型设计和 API 接口实现。
\item \textbf{交叉职责：} 接口契约由前后端共同制定并维护在契约矩阵中；端到端测试由前端主导编写，后端协助提供测试数据和环境。
\end{enumerate}

\subsection{任务管理}
任务管理遵循“需求—>设计—>实现—>验证—>归档”的流程：
\begin{center}
\begin{tikzpicture}[font=\small,
  flowstep/.style={draw,rounded corners,align=center,minimum width=2.3cm,minimum height=.7cm,fill=blue!7},
  arr/.style={-Latex,thick}
]
\node[flowstep] (r) {需求澄清\\编号与验收标准};
\node[flowstep] (d) [right=1.0cm of r] {设计/契约\\页面、接口、模型};
\node[flowstep] (i) [right=1.0cm of d] {实现\\小粒度提交};
\node[flowstep] (t) [below=1.2cm of i] {自动检查\\构建、单测、E2E};
\node[flowstep] (rv) [left=1.0cm of t] {评审/联调\\缺陷与反馈};
\node[flowstep] (rel) [left=1.0cm of rv] {合并/发布\\记录版本};
\draw[arr] (r)--(d);  \draw[arr] (d)--(i);
\draw[arr] (i)--(t);  \draw[arr] (t)--(rv);
\draw[arr] (rv)--(rel);
\draw[arr] (rel.west) .. controls +(-1,0) and +(-1,0) .. (r.west);
\end{tikzpicture}
\vspace{0.2em}
{\small\itshape 图1：任务管理流程}
\end{center}
每个任务由 Issue 或任务卡承载，描述需求编号（如 FR-03）、影响范围、完成定义和风险说明。任务分配基于成员的技能领域和当前负载，优先保证 Must 功能的交付节奏。

\subsection{沟通机制}
团队采用以下沟通机制确保信息同步：
\begin{enumerate}
\item \textbf{接口联调会：} 前后端在开发新功能前共同确认接口字段、枚举值、分页参数和错误语义，形成契约矩阵。接口变更必须经双方确认并更新契约文档，仅在聊天中确认而未更新矩阵不得视为完成。
\item \textbf{代码审查：} 每个合并请求至少由一名非作者成员审查，检查内容包括权限边界、异常路径、代码可读性和测试覆盖。
\item \textbf{变更通知：} 数据库 Schema 变更、接口契约调整或环境配置修改在工作群中同步，确保前后端开发人员及时感知依赖变更。
\end{enumerate}

\section{版本控制规范}
\subsection{Git 仓库管理}
项目分为两个相邻独立仓库：
\begin{enumerate}
\item \textbf{前端仓库（SYSU-Activity-Cloud-Frontend）：} 包含 Vue 3 前端源码、端到端测试、Mock 服务和 CI 工作流。
\item \textbf{后端仓库（Campus-Activity-Information-Platform）：} 包含 Flask 后端源码、pytest 测试、Docker Compose 部署配置和 CI 工作流。
\end{enumerate}
两个仓库在发布时以兼容版本对组成系统基线，记录前端 commit/tag、后端 commit/tag、接口契约版本和数据库 Schema 版本。主分支始终保持可构建状态。

\subsection{分支策略}
\begin{enumerate}
\item \textbf{主分支（main）：} 保持可构建、可部署，禁止直接推送，所有变更通过合并请求引入。
\item \textbf{功能分支（feature/*）：} 从 main 创建，每个分支只解决一个主题。功能完成后通过合并请求合回 main。
\item \textbf{修复分支（fix/*）：} 用于缺陷修复，修复后同时合回 main 和对应的发布分支。
\item \textbf{紧急修复（hotfix/*）：} 从受影响 tag 建立，修复后合回 main。
\end{enumerate}

\subsection{提交规范}
提交信息应简短清晰地描述变更目的，推荐格式为：
\begin{verbatim}
<模块>：<简要描述变更内容>

- <具体改动点 1>
- <具体改动点 2>
\end{verbatim}
示例：\texttt{auth：添加邮箱验证码注册接口}。提交应小而聚焦，避免将多个无关功能的改动混入同一个提交。禁止提交真实密钥、编译产物或本地数据库文件。

\subsection{合并请求与代码审查}
每个合并请求需包含：
\begin{enumerate}
\item 变更说明：关联的需求编号（FR-xx）或缺陷编号，以及变更的简要描述。
\item 验证证据：本地执行 \texttt{npm run verify} 或 \texttt{pytest} 的通过截图或 CI 运行链接。
\item 文档更新：接口契约、配置示例或运维文档若有变更需同步更新。
\end{enumerate}
审查者检查的要点包括：
\begin{itemize}
\item 权限边界是否在服务端校验，前端是否隐藏/禁用越权入口。
\item 异常路径（空数据、网络失败、参数非法）是否被覆盖。
\item 代码可读性：命名是否清晰，逻辑是否过度嵌套，注释是否必要。
\item 测试是否覆盖正向和异常场景，新增代码是否引入未解释的跳过测试。
\end{itemize}
CI 未通过、含真实密钥或破坏主干构建的合并请求禁止合并。紧急缺陷可走快速通道，但必须在修复后补齐测试与复盘记录。

\section{开发规范}
\subsection{代码规范}
\begin{enumerate}
\item \textbf{前端：} 使用 TypeScript 严格模式，禁止使用 \texttt{any} 类型；组件采用组合式 API（Composition API）和 \texttt{<script setup>} 语法；样式使用 CSS 类名，避免内联样式；遵循 Vue 官方风格指南。
\item \textbf{后端：} 遵循 PEP 8 编码规范；视图函数不包含业务逻辑，业务规则沉淀在服务层；数据库操作通过 SQLAlchemy ORM，避免裸 SQL。
\item \textbf{通用：} 不允许硬编码敏感信息（密钥、密码、Token）；不允许提交调试代码（\texttt{console.log}、\texttt{print}、断点）；不允许将编译产物和依赖缓存提交到仓库。
\end{enumerate}

\subsection{目录规范}
前端目录结构按功能分层组织：
\begin{verbatim}
src/
  views/         页面级组件，每个视图一个子目录
  components/    可复用组件
  api/           Axios 接口封装与数据适配
  stores/        Pinia 状态管理
  composables/   可复用逻辑组合式函数
  utils/         纯工具函数
  router/        路由配置
  assets/        静态资源
\end{verbatim}
后端目录结构按架构层次组织：
\begin{verbatim}
app/
  blueprints/    Flask 蓝图（auth、posters、search 等）
  services/      业务服务层
  models/        SQLAlchemy 数据模型
  extensions/    扩展初始化
  utils/         辅助函数
tests/           pytest 测试文件
\end{verbatim}

\subsection{命名规范}
\begin{enumerate}
\item \textbf{前端：} 组件文件使用 PascalCase（\texttt{ActivityDetail.vue}）；普通 TypeScript 文件使用 camelCase（\texttt{activitySort.ts}）；路由路径使用 kebab-case（\texttt{/activity/:id}）。
\item \textbf{后端：} 蓝图使用 snake\_case 模块名；API 端点使用小写复数名词（\texttt{/api/activities}、\texttt{/api/auth/login}）；模型类使用 PascalCase 单数形式（\texttt{User}、\texttt{Poster}、\texttt{Registration}）。
\item \textbf{数据库：} 表名使用 snake\_case 复数形式（\texttt{users}、\texttt{posters}、\texttt{registrations}）；字段名使用 snake\_case（\texttt{created\_at}、\texttt{password\_hash}）。
\end{enumerate}

\subsection{文档规范}
项目文档统一使用 Markdown 格式，存放于 \texttt{docs/} 目录。每份设计文档包含背景、页面结构、接口依赖和验收标准三部分。接口契约矩阵以表格形式记录前端字段名、后端字段名、类型、是否必填和示例值。需求变更时同步更新相关的设计文档、接口契约和测试用例，避免文档与代码脱节。

\section{自动化手段}
\subsection{自动化构建}
前端通过 \texttt{npm run build} 触发 \texttt{vue-tsc -b \&\& vite build}，串行执行 TypeScript 类型检查和 Vite 生产构建，输出可部署的静态资源到 \texttt{dist/} 目录。后端通过 Dockerfile 定义构建过程，安装 Python 依赖并打包 Flask 应用。

\subsection{自动化检查}
\begin{enumerate}
\item \textbf{类型检查：} \texttt{vue-tsc} 在构建时对 TypeScript 代码执行严格类型检查，捕获接口字段不匹配和类型错误。
\item \textbf{代码质量：} 后端 CI 计划集成 pylint 或 flake8 进行静态检查，当前阶段以前端构建时的类型检查为主要质量门禁。
\item \textbf{安全门禁：} 合并请求审查中对敏感信息泄露、硬编码凭据和 SSRF 风险进行人工检查。
\end{enumerate}

\subsection{自动化测试}
\begin{enumerate}
\item \textbf{单元测试：} 前端使用 Vitest + \texttt{@vue/test-utils}，覆盖认证 Store、活动排序/质量工具、日历卡片和活动详情等组件。后端使用 pytest + 内存 SQLite，覆盖认证、活动、审计、数据源、知识、搜索、质量、去重和安全等模块。
\item \textbf{端到端测试：} 使用 Playwright 驱动 Chromium 浏览器，覆盖访客浏览、登录注册、活动搜索等关键用户旅程。
\item \textbf{聚合验证：} \texttt{npm run verify} 串行执行构建、单元测试和端到端测试，作为提交前的完整性检查命令。
\end{enumerate}

\subsection{自动化打包}
前端构建生成静态文件（HTML/CSS/JS），可直接部署到 Web 服务器或反向代理后端。后端通过 Docker 构建不可变镜像，由 Compose 编排启动。CI 工作流在 push 和 PR 时自动触发，在 Ubuntu 环境中执行 npm ci、浏览器安装、构建和测试。

\section{项目管理工具}
项目使用以下工具支撑工程化实践：
\begin{enumerate}
\item \textbf{GitHub：} 代码托管、版本控制和协作平台。前端和后端分属两个独立仓库，通过 Issue 管理任务和缺陷。
\item \textbf{GitHub Actions：} 持续集成平台。前端仓库已配置 CI 工作流，在 push 和 pull request 上自动执行安装、构建、单元测试和端到端测试。
\item \textbf{LaTeX：} 课程报告编写工具。六份设计文档均通过 XeLaTeX 编译生成 PDF。
\item \textbf{Docker Compose：} 部署编排工具。定义 API、PostgreSQL、Redis、Celery Worker/Beat 的容器化拓扑，单命令启动完整后端服务。
\item \textbf{Mock 服务：} 基于 Node.js 的轻量级 API Mock，供前端在无后端环境时独立开发和联调。
\end{enumerate}

\section{工程化实践效果}
当前工程化实践的落地情况如下：
\begin{tabularx}{\textwidth}{p{2.8cm}p{3.5cm}Y}
\toprule 实践 & 仓库证据 & 当前状态 \\\midrule
前端构建与类型检查 & \texttt{package.json} 的 build/verify 脚本 & 已配置并集成到 CI \\
前端单元测试 & Vitest 覆盖 5 个组件/工具模块 & 已配置并集成到 CI \\
前端 E2E 测试 & Playwright 覆盖公开访问和认证流 & 已配置并集成到 CI \\
前端 CI & \texttt{.github/workflows/ci.yml}（push/PR 触发） & 已实现自动化 \\
前端设计文档 & \texttt{docs/frontend-design/*.md}（8 份设计文档） & 已形成文档资产 \\
接口契约矩阵 & \texttt{docs/frontend-plan/BACKEND\_CONTRACT\_MATRIX.md} & 已建立并在迭代中更新 \\
后端测试资产 & \texttt{backend/tests/}（20+ 测试文件） & 测试代码存在，CI 待补充 \\
后端部署配置 & Dockerfile、Compose、Gunicorn、\texttt{.env.example} & 已配置作为部署起点 \\
\bottomrule
\end{tabularx}

当前已建立以自动化前端 CI 和文档化为核心的工程基础，后端 CI、静态检查门禁、覆盖率阈值和契约自动化测试被明确列为下一步改进方向。改进优先级为：
\begin{enumerate}
\item 新建后端 CI 工作流：安装开发依赖后执行 pytest、静态检查和覆盖率报告。
\item 将契约矩阵纳入自动化契约测试，阻止字段和分页回归。
\item 以受保护主分支、必需评审和必需 CI 状态实施分支保护策略。
\item 发布流水线增加镜像版本标签、健康检查、迁移执行和回滚步骤。
\item 用 Issue 看板统计缺陷和交付周期数据，在迭代回顾中驱动过程改进。
\end{enumerate}
以上改进均为计划项，在当前报告中不做为已完成能力表述。

\end{document}
